\documentclass[english,twoside, openright, hidelinks, 11 pt, class=report,crop=false]{standalone}
\input{../../preamb}
\input{../bm}

\begin{document}
\subsection*{Kapittel \ref{Geo}}
\opr{opggeoviscos} 
\abch{
	\item \selos $\sin 45^\circ=\frac{\sqrt{2}}{2}$.
	\item \selos
	\item \selos
	\item $\sin 60^\circ=\frac{\sqrt{3}}{2}$ og $\cos 60^\circ=\frac{1}{2}$
	\item $\tan 45^\circ=1$, $\tan 30^\circ=\frac{1}{\sqrt{3}}$, $\tan 60^\circ=\sqrt{3}$
}

\opr{opggeouminv}
\selos

\opr{opggeovistan}
\selos

\opr{opggeoquatrants}
\selos

\opr{opggeofindsincostan}
\abch{
	\item $\sin 150^\circ=\frac{1}{2}$, $\cos 150^\circ=-\frac{\sqrt{3}}{2} $, $\tan 150^\circ=-\frac{1}{\sqrt{3}}$
	\item $\sin 135^\circ=\frac{\sqrt{2}}{2}$, $\cos 135^\circ=-\frac{\sqrt{2}}{2}$, $\tan 135^\circ=-1$
	\item $\sin 120^\circ=\frac{\sqrt{3}}{2}$, $\cos 120^\circ=-\frac{1}{2}$, $\tan 120^\circ=-\sqrt{3}$
}

\opr{opggeoperf}
\abch{
\item $v=60^\circ$
\item $v=10^\circ$
}


\opr{opggeoarsetn}
\abch{
	\item $\frac{35\sqrt{3}}{4}$
	\item $\frac{3(\sqrt{5}-1)}{2	}$
	\item $ \frac{3+\sqrt{3}}{2} $
}

\opr{opggeolikeperf}
Vinklene har samme sentralvinkel, og må dermed ha lik størrelse.


\opr{opggeotalesTM} \vs
\abch{
	\item \selos
	\item $\angle C=90^\circ \Leftarrow \text{$ AB $ er en diameter i den omskrevne sirkelen til }\triangle ABC$
}

\opr{opggeotanchord}
\selos

\opr{opggeoabcr}
\selos

\opr{opggeo2r}
\abch{
	\item \selos
	\item \selos
	\item \selos
}

\opr{1tv21d1opg2}
6

\opr{1tv25d1opg5} 
\abch{
	\item 15
	\item 7
}

\opr{opggeoeqar}
$4\sqrt{3}$

\opr{1TV23D1opg1}
\abch{
	\item \selos
	\item \selos
	\item \selos
	\item \selos
}

\opr{opggeoarealsetn}
\abch{
	\item \selos
	\item \selos
}

\opr{1th21d1opg2}
$AC=8$

\opr{1tv26d1opg6}
Trekanten er rettvinklet og likebeint.

\opr{1TH21D1opg8}
$BC=\sqrt{28}$

\grubr{opggeoslope}
\selos

\opr{1tv26d1opg7}
\abch{
\item $10\sqrt{3}$
\item $18+10\sqrt{3}$
}

\grubr{1tv24d2opg3}
$AC=2$, $BC=2\sqrt{21}$

%\grubr{1TH21D1opg9} 
%$\dfrac{24}{5}$

\grubr{1th22d2opg3}
\abch{
	\item 2
	\item $6\sqrt{2}$
}


\grubr{1tv23d2opg5}
4

\grubr{1th22d2opg4}
\abch{
	\item Den ene sidelengden er $14$. Vinkelen mellom sidene med lengder $x$ og $16$ er $60^\circ$.
	\item $6$, $14$ og $16$ eller $10$, $14$ og $16$
	\item $8\sqrt{3}$
}

\grubr{opggeo1575}
$\cos 15^\circ=\frac{1+\sqrt{3}}{2\sqrt{2}}$, $\sin 75^\circ=\frac{1+\sqrt{3}}{2\sqrt{2}}$, $\cos 75^\circ=\frac{\sqrt{3}-1}{2\sqrt{2}}$ og $\tan 15^\circ=2-\sqrt{3}$

\grubr{1tv22d2opg3} 
$12+4\sqrt{3}$


\grubr{1TV21D1opg2} $\sin 160^\circ < \log 1 < \sin 60^\circ < \left(\frac{3}{4}\right)^{-1}$


\grubr{t1h23d1opg1}
\abch{
	\item \selos
	\item \selos	
}

\grubr{1tv22d1opg3} 
\abch{
	\item Nei.
	\item Ja.
	\item Nei
}

\grubr{t1h23d1opg4}
Trekanten med vinkel $32^\circ$ markert.

\grubr{1tv25d2opg3}
\abch{
	\item $2\sqrt{10}$
	\item $24\sqrt{5}(\sqrt{3}-1)$
}

\grubr{opggeo180minv}
\selos

\grubr{opggeokordteo}
\selos

\grubr{opggeobeviscossetn}
\selos

\grubr{opggeovinkar}
$r^2\left(\frac{\sin(2v)}{2}+\frac{\pi v}{180^\circ}\right)$

\grubr{opggeolimsin0overx}
\selos

\grubr{opggeoR}
\selos

\grubr{opggeoviscosuv}
\selos

\grubr{opggeotalesTMconv}
\selos

\grubr{opgeo45vink}
$45^\circ$

\grubr{opggeotangentsq}
\selos

\grubr{opggeosin18}
\selos


\end{document}


